\documentclass[11pt,a4paper]{article}

\usepackage[utf8]{inputenc}
\usepackage[T1]{fontenc}
\usepackage[spanish,es-nodecimaldot]{babel}
\usepackage{lmodern}
\usepackage{geometry}
\usepackage{xcolor}
\usepackage{hyperref}
\usepackage{xurl}
\usepackage{enumitem}
\usepackage{booktabs}
\usepackage{tabularx}
\usepackage{array}
\usepackage{listings}
\usepackage{fancyhdr}

\geometry{top=2cm,bottom=2cm,left=2.1cm,right=2.1cm}
\definecolor{AzulSipeip}{HTML}{0B2A45}
\definecolor{CelesteSipeip}{HTML}{146C94}
\definecolor{FondoCodigo}{HTML}{F5F7F9}
\definecolor{BordeCodigo}{HTML}{D7DEE5}
\definecolor{VerdeResultado}{HTML}{176B3A}
\definecolor{RojoAtencion}{HTML}{9A3412}

\hypersetup{
  colorlinks=true,
  linkcolor=AzulSipeip,
  urlcolor=CelesteSipeip,
  pdftitle={Guía práctica de pruebas de SIPeIP en Postman}
}

\pagestyle{fancy}
\fancyhf{}
\fancyhead[L]{\textcolor{AzulSipeip}{SIPeIP}}
\fancyhead[R]{\small Pruebas manuales en Postman}
\fancyfoot[C]{\thepage}
\renewcommand{\headrulewidth}{0.4pt}

\setlength{\parindent}{0pt}
\setlength{\parskip}{0.55em}
\setlist[enumerate]{leftmargin=1.7em,itemsep=0.3em,topsep=0.3em}
\setlist[itemize]{leftmargin=1.5em,itemsep=0.25em,topsep=0.3em}
\newcolumntype{Y}{>{\raggedright\arraybackslash}X}

\lstdefinestyle{postman}{
  backgroundcolor=\color{FondoCodigo},
  basicstyle=\ttfamily\small,
  breaklines=true,
  columns=fullflexible,
  frame=single,
  rulecolor=\color{BordeCodigo},
  framerule=0.5pt,
  showstringspaces=false,
  xleftmargin=0.2em,
  xrightmargin=0.2em,
  aboveskip=0.4em,
  belowskip=0.55em,
  literate=
    {á}{{\'a}}1 {é}{{\'e}}1 {í}{{\'i}}1 {ó}{{\'o}}1 {ú}{{\'u}}1
    {Á}{{\'A}}1 {É}{{\'E}}1 {Í}{{\'I}}1 {Ó}{{\'O}}1 {Ú}{{\'U}}1
    {ñ}{{\~n}}1 {Ñ}{{\~N}}1 {ü}{{\"u}}1
}
\lstset{style=postman}

\newcommand{\esperado}[1]{\textcolor{VerdeResultado}{\textbf{Resultado esperado:}} #1}
\newcommand{\atencion}[1]{\textcolor{RojoAtencion}{\textbf{Importante:}} #1}
\newcommand{\metodo}[1]{\textbf{Método:} \texttt{#1}}

\begin{document}

{\LARGE\bfseries\color{AzulSipeip} Guía práctica para probar SIPeIP en Postman}

Esta guía empieza directamente con las pruebas. El backend debe estar disponible en \url{http://localhost:8000} y los seis usuarios ya deben existir. No es necesario configurar Bearer Token ni escribir cookies manualmente.

En los datos demostrativos, \textbf{SNP} corresponde a la Secretaría Nacional de Planificación y \textbf{MDI} es una institución pública externa ficticia usada únicamente para probar el aislamiento entre entidades.

\atencion{Las pruebas funcionales crean registros reales en la base local, aprueban un plan demostrativo y agregan avances. Los nombres y códigos GUIA permiten reconocer esos registros. Para repetir una creación use otro sufijo, por ejemplo GUIA-02.}

\section{Lo único que debe saber antes de enviar solicitudes}

\subsection{Postman conservará la sesión}

SIPeIP inicia sesión mediante cookies. Postman guarda esas cookies automáticamente después de un login correcto siempre que todas las URLs usen \path{localhost}. No cambie algunas URLs a \path{127.0.0.1}.

En la pestaña \textbf{Authorization} de cada solicitud seleccione \textbf{No Auth}. No agregue \path{Authorization}, \path{Bearer Token}, \path{Cookie} ni \path{sessionid} por su cuenta.

\subsection{Qué es el token CSRF y dónde se coloca}

Antes de enviar un \path{POST}, \path{PATCH}, \path{PUT} o \path{DELETE}, SIPeIP exige un código temporal llamado CSRF. Ese código demuestra que la solicitud pertenece a la misma sesión de Postman.

Para obtenerlo:

\begin{enumerate}
  \item Cree una solicitud nueva en Postman.
  \item Seleccione el método \textbf{GET}.
  \item Pegue esta URL y pulse \textbf{Send}:
\end{enumerate}

\begin{lstlisting}
http://localhost:8000/api/auth/csrf/
\end{lstlisting}

La respuesta será parecida a esta:

\begin{lstlisting}
{
  "detail": "Token CSRF preparado correctamente.",
  "csrf_token": "valor-largo-generado-por-django"
}
\end{lstlisting}

Copie únicamente el contenido de \path{csrf_token}, sin las comillas.

Cuando esta guía indique un \path{POST}, \path{PATCH} o \path{DELETE}, abra la pestaña \textbf{Headers} de esa solicitud y agregue exactamente estas dos filas:

\begin{center}
\begin{tabularx}{0.9\textwidth}{Y Y}
\toprule
\textbf{Key} & \textbf{Value} \\
\midrule
\path{Content-Type} & \path{application/json} \\
\path{X-CSRFToken} & Pegue aquí el valor de \path{csrf_token} \\
\bottomrule
\end{tabularx}
\end{center}

Para colocar un JSON, abra \textbf{Body}, seleccione \textbf{raw} y luego seleccione \textbf{JSON} en la lista que aparece a la derecha.

\atencion{Django cambia el token CSRF cuando el login es correcto. Por eso, inmediatamente después de iniciar sesión debe volver a ejecutar \path{GET /api/auth/csrf/} y copiar el nuevo token.}

\subsection{Cómo cambiar de usuario durante las pruebas}

Cada vez que la guía pida usar otro perfil:

\begin{enumerate}
  \item Si hay una sesión abierta y no recuerda el token actual, ejecute primero \path{GET http://localhost:8000/api/auth/csrf/} y cópielo.
  \item Cree una solicitud \path{POST} a \url{http://localhost:8000/api/auth/logout/}.
  \item En \textbf{Headers} agregue \path{Content-Type = application/json} y \path{X-CSRFToken = token actual}.
  \item En \textbf{Body $>$ raw $>$ JSON} escriba \texttt{\{\}} y pulse \textbf{Send}.
\end{enumerate}

El logout correcto responde status \textbf{200}. Si responde 401 porque la sesión ya había expirado, puede continuar con el nuevo login.

Para iniciar la nueva sesión:

\begin{enumerate}
  \item Ejecute \path{GET http://localhost:8000/api/auth/csrf/} y copie el token recibido.
  \item Cree una solicitud \textbf{POST} a \url{http://localhost:8000/api/auth/login/}.
  \item En \textbf{Headers} agregue \path{Content-Type = application/json} y \path{X-CSRFToken = token que acaba de copiar}.
  \item En \textbf{Body $>$ raw $>$ JSON} escriba:
\end{enumerate}

\begin{lstlisting}
{
  "username": "<USUARIO>",
  "password": "<CONTRASENA>"
}
\end{lstlisting}

Reemplace \path{<USUARIO>}, pulse \textbf{Send} y confirme status \textbf{200}. Después ejecute de nuevo \path{GET http://localhost:8000/api/auth/csrf/} y copie el token nuevo que usará en las demás solicitudes.

Los usuarios disponibles son \path{administrador}, \path{planificador}, \path{supervisor}, \path{externo}, \path{auditor} y \path{superadministrador}. En los ejemplos, reemplace \path{<CONTRASENA>} por la contraseña que usó al crearlos.

\subsection{Cómo usar los valores entre signos menor y mayor}

Un texto como \path{<ID_PLAN>} es un marcador. Debe reemplazarlo por el número real que aparece en la respuesta de Postman. Por ejemplo, si un plan devuelve \path{"id": 7}, esta URL:

\begin{lstlisting}
http://localhost:8000/api/planes/<ID_PLAN>/
\end{lstlisting}

se convierte en:

\begin{lstlisting}
http://localhost:8000/api/planes/7/
\end{lstlisting}

No envíe los signos \path{< >} literalmente.

Cuando una URL incluya un signo \path{?}, lo que aparece después son parámetros de consulta. Puede pegar la URL completa tal como está escrita; Postman mostrará automáticamente esas parejas en la pestaña \textbf{Params}. No vuelva a agregarlas una segunda vez.

\section{Prueba 1: autenticación, sesión y CSRF}

\subsection{AUTH-01 --- Intentar entrar sin iniciar sesión}

\begin{enumerate}
  \item Asegúrese de no tener una sesión abierta. Si tiene una, use el logout explicado arriba. Si no puede cerrarla, abra \textbf{Cookies} en Postman, seleccione \path{localhost} y elimine únicamente la cookie \path{sessionid}.
  \item Cree una solicitud y seleccione el método \textbf{GET}.
  \item Pegue la URL y pulse \textbf{Send}:
\end{enumerate}

\begin{lstlisting}
http://localhost:8000/api/auth/me/
\end{lstlisting}

No seleccione Body y no agregue filas manuales en \textbf{Headers}.

\esperado{Status \textbf{401}. El JSON debe incluir \path{"code": "not_authenticated"}. Esto comprueba que un recurso protegido no puede consultarse sin sesión.}

\subsection{AUTH-02 --- Comprobar que CSRF es obligatorio}

\begin{enumerate}
  \item Ejecute \path{GET /api/auth/csrf/} para que Postman reciba la cookie CSRF.
  \item Cree una solicitud \textbf{POST} con esta URL:
\end{enumerate}

\begin{lstlisting}
http://localhost:8000/api/auth/login/
\end{lstlisting}

En \textbf{Headers} agregue solamente esta fila y omita intencionalmente \path{X-CSRFToken}:

\begin{center}
\begin{tabular}{ll}
\toprule
\textbf{Key} & \textbf{Value} \\
\midrule
\path{Content-Type} & \path{application/json} \\
\bottomrule
\end{tabular}
\end{center}

En \textbf{Body $>$ raw $>$ JSON} pegue:

\begin{lstlisting}
{
  "username": "planificador",
  "password": "<CONTRASENA>"
}
\end{lstlisting}

Pulse \textbf{Send}.

\esperado{Status \textbf{403} y \path{"code": "csrf_failed"}. La sesión no debe iniciarse.}

\subsection{AUTH-03 --- Probar una contraseña incorrecta}

\begin{enumerate}
  \item Vuelva a ejecutar \path{GET http://localhost:8000/api/auth/csrf/} y copie el token.
  \item Cree una solicitud \textbf{POST} a \path{/api/auth/login/}.
  \item En \textbf{Headers} agregue las dos filas siguientes:
\end{enumerate}

\begin{center}
\begin{tabular}{ll}
\toprule
\textbf{Key} & \textbf{Value} \\
\midrule
\path{Content-Type} & \path{application/json} \\
\path{X-CSRFToken} & El token que acaba de copiar \\
\bottomrule
\end{tabular}
\end{center}

En \textbf{Body $>$ raw $>$ JSON} pegue:

\begin{lstlisting}
{
  "username": "planificador",
  "password": "valor-incorrecto-controlado"
}
\end{lstlisting}

\esperado{Status \textbf{401}. El mensaje no debe revelar si falló el nombre de usuario o la contraseña.}

\subsection{AUTH-04 --- Iniciar sesión correctamente}

\begin{enumerate}
  \item Cree una solicitud \textbf{POST} a \url{http://localhost:8000/api/auth/login/}.
  \item En \textbf{Headers} agregue \path{Content-Type = application/json} y \path{X-CSRFToken = token obtenido en AUTH-03}.
  \item En \textbf{Body $>$ raw $>$ JSON} pegue:
\end{enumerate}

\begin{lstlisting}
{
  "username": "planificador",
  "password": "<CONTRASENA>"
}
\end{lstlisting}

Pulse \textbf{Send}.

\esperado{Status \textbf{200}. La respuesta identifica a \path{planificador}, muestra su rol, entidad, permisos y tiempo de expiración. La contraseña y la cookie de sesión no aparecen en el JSON.}

Ahora haga obligatoriamente lo siguiente:

\begin{enumerate}
  \item Ejecute nuevamente \path{GET http://localhost:8000/api/auth/csrf/}.
  \item Copie el nuevo \path{csrf_token}. Este será el token que usará en las siguientes solicitudes.
  \item Ejecute \path{GET http://localhost:8000/api/auth/me/}.
\end{enumerate}

\esperado{La consulta \path{/me/} devuelve status \textbf{200} e identifica al usuario \path{planificador}.}

\subsection{AUTH-05 --- Renovar la sesión}

\begin{enumerate}
  \item Cree una solicitud \textbf{POST} a \url{http://localhost:8000/api/auth/refresh/}.
  \item En \textbf{Headers} agregue \path{Content-Type = application/json} y \path{X-CSRFToken = token obtenido después del login}.
  \item En \textbf{Body $>$ raw $>$ JSON} escriba \texttt{\{\}}.
  \item Pulse \textbf{Send}.
\end{enumerate}

\esperado{Status \textbf{200}. La respuesta indica que la sesión fue renovada y vuelve a informar \path{expira_en_segundos}.}

\subsection{AUTH-06 --- Consultar el dashboard del perfil}

Cree una solicitud \textbf{GET}; no seleccione Body ni agregue filas manuales en \textbf{Headers}. Use esta URL:

\begin{lstlisting}
http://localhost:8000/api/dashboard/
\end{lstlisting}

\esperado{Status \textbf{200}. La respuesta contiene \path{alcance}, \path{entidad} y una lista de \path{widgets} permitidos para \path{planificador}.}

\section{Prueba 2: localizar los IDs que se usarán después}

Cambie al usuario \path{superadministrador} siguiendo la sección ``Cómo cambiar de usuario''. Luego cree solicitudes \textbf{GET} para las siguientes URLs. No seleccione Body ni agregue filas manuales en \textbf{Headers}.

\begin{lstlisting}
http://localhost:8000/api/roles/?search=Usuario%20Externo
http://localhost:8000/api/roles/?search=Superadministrador
http://localhost:8000/api/configuracion/entidades/?search=MDI
http://localhost:8000/api/configuracion/entidades/?search=SNP
http://localhost:8000/api/configuracion/unidades/?search=UADM
http://localhost:8000/api/usuarios/?search=planificador
http://localhost:8000/api/usuarios/?search=supervisor
http://localhost:8000/api/ejes-pnd/?estado=ACTIVO
http://localhost:8000/api/ods/?estado=ACTIVO
http://localhost:8000/api/proyectos/?search=PIR-001
http://localhost:8000/api/hitos-proyectos/?search=PIR-001
\end{lstlisting}

En cada respuesta busque el campo \path{id} del registro correcto y anótelo temporalmente:

\begin{itemize}
  \item ID del rol Usuario Externo;
  \item ID del rol Superadministrador técnico;
  \item ID de las entidades MDI y SNP;
  \item ID de la unidad UADM de MDI;
  \item ID de los usuarios planificador y supervisor;
  \item ID del eje EJE-REF-01;
  \item ID del ODS número 16;
  \item ID del proyecto PIR-001 y de su hito.
\end{itemize}

Los registros \path{EJE-REF-01}, \path{PND-REF-01} y ODS 16 son referencias demostrativas y configurables de esta base local; la guía no los presenta como un catálogo oficial completo.

Con el ID del eje PND, ejecute también:

\begin{lstlisting}
http://localhost:8000/api/objetivos-pnd/?eje=<ID_EJE_PND>&estado=ACTIVO
\end{lstlisting}

Guarde el ID del objetivo PND con código \path{PND-REF-01}.

\esperado{Todas las solicitudes responden status \textbf{200}. Los IDs pueden ser distintos en cada base, por eso deben copiarse de las respuestas reales.}

\section{Prueba 3: permisos por rol y protección técnica}

SIPeIP usa control de acceso basado en roles, conocido como \textbf{RBAC}. En la práctica, cada rol agrupa permisos concretos y la API vuelve a comprobarlos en cada operación.

\subsection{RBAC-01 --- Consultar el catálogo real de permisos}

Con la sesión de \path{superadministrador}:

\begin{enumerate}
  \item Cree una solicitud \textbf{GET}.
  \item Pegue esta URL y pulse \textbf{Send}:
\end{enumerate}

\begin{lstlisting}
http://localhost:8000/api/roles/catalogo-permisos/
\end{lstlisting}

No seleccione Body y no agregue filas manuales en \textbf{Headers}.

\esperado{Status \textbf{200} y una lista llamada \path{permisos}. Esta lista proviene del backend.}

\subsection{RBAC-02 --- Impedir que un administrador modifique el rol técnico}

Cambie al usuario \path{administrador} y obtenga el nuevo token CSRF después del login.

\begin{enumerate}
  \item Cree una solicitud y seleccione \textbf{PATCH}.
  \item Use esta URL reemplazando el ID:
\end{enumerate}

\begin{lstlisting}
http://localhost:8000/api/roles/<ID_ROL_SUPERADMIN>/
\end{lstlisting}

En \textbf{Headers} agregue:

\begin{center}
\begin{tabular}{ll}
\toprule
\textbf{Key} & \textbf{Value} \\
\midrule
\path{Content-Type} & \path{application/json} \\
\path{X-CSRFToken} & El token obtenido después del login \\
\bottomrule
\end{tabular}
\end{center}

En \textbf{Body $>$ raw $>$ JSON} pegue:

\begin{lstlisting}
{
  "descripcion": "Intento controlado desde Postman"
}
\end{lstlisting}

Pulse \textbf{Send}.

\esperado{Status \textbf{403}. El Administrador del Sistema no puede modificar el rol técnico superior.}

\section{Prueba 4: crear, validar, bloquear y eliminar un usuario}

Cambie a \path{superadministrador} y obtenga el token CSRF posterior al login.

\subsection{USR-01 --- Crear un usuario temporal}

\begin{enumerate}
  \item Cree una solicitud \textbf{POST}.
  \item Use la URL:
\end{enumerate}

\begin{lstlisting}
http://localhost:8000/api/usuarios/
\end{lstlisting}

En \textbf{Headers} agregue:

\begin{center}
\begin{tabular}{ll}
\toprule
\textbf{Key} & \textbf{Value} \\
\midrule
\path{Content-Type} & \path{application/json} \\
\path{X-CSRFToken} & El token actual \\
\bottomrule
\end{tabular}
\end{center}

En \textbf{Body $>$ raw $>$ JSON} pegue el siguiente contenido después de reemplazar los IDs y elegir una contraseña temporal de al menos ocho caracteres, que no sea común ni se parezca al nombre del usuario:

\begin{lstlisting}
{
  "username": "guia.usuario.01",
  "email": "guia.usuario.01@example.test",
  "first_name": "Usuario",
  "last_name": "Guia",
  "password": "<CONTRASENA_TEMPORAL>",
  "rol": <ID_ROL_EXTERNO>,
  "telefono": "0999999999",
  "entidad": <ID_ENTIDAD_MDI>,
  "unidad_organizacional": <ID_UNIDAD_UADM>
}
\end{lstlisting}

\esperado{Status \textbf{201}. Copie el \path{id} del usuario creado. La respuesta debe mostrar estado ACTIVO, pero no debe devolver la contraseña.}

\subsection{USR-02 --- Comprobar una validación}

\begin{enumerate}
  \item Cree una solicitud \textbf{PATCH}.
  \item Use \path{http://localhost:8000/api/usuarios/<ID_USUARIO_CREADO>/}.
  \item En \textbf{Headers} agregue \path{Content-Type = application/json} y \path{X-CSRFToken = token actual}.
  \item En \textbf{Body $>$ raw $>$ JSON} pegue:
\end{enumerate}

\begin{lstlisting}
{
  "telefono": "12345"
}
\end{lstlisting}

\esperado{Status \textbf{400} y un mensaje en español indicando que el teléfono debe tener diez dígitos y comenzar con 09.}

Sin cambiar la URL ni los headers, reemplace el Body por:

\begin{lstlisting}
{
  "telefono": "0987654321"
}
\end{lstlisting}

\esperado{Status \textbf{200}. El teléfono queda actualizado.}

\subsection{USR-03 --- Bloquear y volver a activar el usuario}

Primero bloquee el usuario:

\begin{enumerate}
  \item Seleccione el método \textbf{POST}.
  \item Use \path{http://localhost:8000/api/usuarios/<ID_USUARIO_CREADO>/bloquear/}.
  \item En \textbf{Headers} agregue \path{Content-Type = application/json} y \path{X-CSRFToken = token actual}.
  \item En \textbf{Body $>$ raw $>$ JSON} escriba \texttt{\{\}}.
  \item Pulse \textbf{Send}.
\end{enumerate}

\esperado{Status \textbf{200}, estado BLOQUEADO e \path{is_active: false}.}

Después cambie únicamente la URL a:

\begin{lstlisting}
http://localhost:8000/api/usuarios/<ID_USUARIO_CREADO>/activar/
\end{lstlisting}

Mantenga el método POST, los dos headers y el Body \texttt{\{\}}.

\esperado{Status \textbf{200}, estado ACTIVO e \path{is_active: true}.}

\subsection{USR-04 --- Eliminar el usuario temporal}

No inicie sesión con ese usuario ni lo asigne a un plan antes de esta prueba.

\begin{enumerate}
  \item Cree una solicitud \textbf{DELETE}.
  \item Use \path{http://localhost:8000/api/usuarios/<ID_USUARIO_CREADO>/}.
  \item En \textbf{Headers} agregue \path{Content-Type = application/json} y \path{X-CSRFToken = token actual}.
  \item No agregue Body.
  \item Pulse \textbf{Send}.
\end{enumerate}

\esperado{Status \textbf{204}. Si el usuario adquirió relaciones trazables, la API debe responder 409 y conservarlo.}

\section{Prueba 5: comprobar que un Usuario Externo solo ve su institución}

Cambie al usuario \path{externo}.

\subsection{SCOPE-01 --- Listar entidades visibles}

\begin{enumerate}
  \item Cree una solicitud \textbf{GET}.
  \item Use \url{http://localhost:8000/api/configuracion/entidades/}.
  \item No seleccione Body y no agregue filas manuales en \textbf{Headers}.
  \item Pulse \textbf{Send}.
\end{enumerate}

\esperado{Status \textbf{200}. La respuesta contiene MDI y no contiene SNP.}

\subsection{SCOPE-02 --- Intentar consultar directamente la entidad SNP}

Use una solicitud \textbf{GET} con:

\begin{lstlisting}
http://localhost:8000/api/configuracion/entidades/<ID_ENTIDAD_SNP>/
\end{lstlisting}

No seleccione Body y no agregue filas manuales en \textbf{Headers}.

\esperado{Status \textbf{404}. Aunque SNP exista, la API usa esta respuesta para no revelar registros de otra institución al Usuario Externo.}

\subsection{SCOPE-03 --- Intentar crear una entidad sin permiso}

\begin{enumerate}
  \item Cree una solicitud \textbf{POST} a \url{http://localhost:8000/api/configuracion/entidades/}.
  \item En \textbf{Headers} agregue \path{Content-Type = application/json} y \path{X-CSRFToken = token actual}.
  \item En \textbf{Body $>$ raw $>$ JSON} pegue:
\end{enumerate}

\begin{lstlisting}
{
  "codigo_oficial": "GUIA-NO-AUT",
  "nombre": "Entidad que no debe crearse",
  "subsector": "Prueba controlada",
  "nivel_gobierno": "Nacional"
}
\end{lstlisting}

\esperado{Status \textbf{403}. El Usuario Externo puede consultar su configuración, pero no administrarla.}

\section{Prueba 6: objetivo estratégico y alineación PND/ODS}

En esta sección, \textbf{PND} significa Plan Nacional de Desarrollo y \textbf{ODS} significa Objetivos de Desarrollo Sostenible.

Cambie al usuario \path{planificador}.

\subsection{OBJ-01 --- Crear un objetivo estratégico de MDI}

\begin{enumerate}
  \item Cree una solicitud \textbf{POST} a \url{http://localhost:8000/api/objetivos-estrategicos/}.
  \item En \textbf{Headers} agregue \path{Content-Type = application/json} y \path{X-CSRFToken = token actual}.
  \item En \textbf{Body $>$ raw $>$ JSON} pegue:
\end{enumerate}

\begin{lstlisting}
{
  "entidad": <ID_ENTIDAD_MDI>,
  "codigo": "OEI-GUIA-01",
  "nombre": "Fortalecer la gestion institucional",
  "descripcion": "Objetivo creado durante la guia de pruebas."
}
\end{lstlisting}

\esperado{Status \textbf{201}. Copie el ID del objetivo. Debe aparecer en estado ACTIVO.}

Si el código \path{OEI-GUIA-01} ya existe porque repitió la guía, use \path{OEI-GUIA-02}.

\subsection{ALI-01 --- Relacionar el objetivo con PND y ODS}

\begin{enumerate}
  \item Cree una solicitud \textbf{POST} a \url{http://localhost:8000/api/alineaciones/}.
  \item En \textbf{Headers} agregue \path{Content-Type = application/json} y \path{X-CSRFToken = token actual}.
  \item En \textbf{Body $>$ raw $>$ JSON} pegue, reemplazando los tres IDs:
\end{enumerate}

\begin{lstlisting}
{
  "objetivo_estrategico": <ID_OBJETIVO_CREADO>,
  "objetivo_pnd": <ID_OBJETIVO_PND>,
  "ods": <ID_ODS_16>,
  "justificacion": "La mejora institucional contribuye al objetivo PND y al ODS seleccionados."
}
\end{lstlisting}

\esperado{Status \textbf{201}, estado BORRADOR y \path{usuario_creador_detalle.username: "planificador"}. Copie el ID de la alineación.}

\subsection{ALI-02 --- Comprobar quién puede validar la alineación}

Como \path{planificador}:

\begin{enumerate}
  \item Cree una solicitud \textbf{POST} a \path{http://localhost:8000/api/alineaciones/<ID_ALINEACION>/validar/}.
  \item En \textbf{Headers} agregue \path{Content-Type = application/json} y \path{X-CSRFToken = token actual}.
  \item En \textbf{Body $>$ raw $>$ JSON} escriba \texttt{\{\}}.
  \item Pulse \textbf{Send}.
\end{enumerate}

\esperado{Status \textbf{403}. Crear una alineación no significa poder validarla.}

Cambie a \path{superadministrador}, obtenga su token CSRF posterior al login y repita exactamente la misma solicitud con el nuevo token.

\esperado{Status \textbf{200}, estado VALIDADA y \path{usuario_validador_detalle.username: "superadministrador"}. Si vuelve a pulsar Send, debe obtener 409 porque una alineación validada ya no está en borrador.}

\section{Prueba 7: flujo completo Plan, Meta, Indicador y Avance}

Cambie al usuario \path{planificador}. Use el objetivo estratégico creado en OBJ-01.

\subsection{PLAN-01 --- Crear un plan}

\begin{enumerate}
  \item Cree una solicitud \textbf{POST} a \url{http://localhost:8000/api/planes/}.
  \item En \textbf{Headers} agregue \path{Content-Type = application/json} y \path{X-CSRFToken = token actual}.
  \item En \textbf{Body $>$ raw $>$ JSON} pegue:
\end{enumerate}

\begin{lstlisting}
{
  "nombre": "Plan Postman GUIA-01",
  "descripcion": "Plan creado para comprobar el flujo completo.",
  "entidad": <ID_ENTIDAD_MDI>,
  "periodo_inicio": "2026-01-01",
  "periodo_fin": "2026-12-31",
  "responsable": <ID_USUARIO_PLANIFICADOR>
}
\end{lstlisting}

\esperado{Status \textbf{201} y estado BORRADOR. Copie el ID del plan.}

Si el nombre ya existe, cambie \path{GUIA-01} por otro número.

\subsection{PLAN-02 --- Probar la validación de fechas}

\begin{enumerate}
  \item Cree una solicitud \textbf{PATCH} a \path{http://localhost:8000/api/planes/<ID_PLAN>/}.
  \item En \textbf{Headers} agregue \path{Content-Type = application/json} y \path{X-CSRFToken = token actual}.
  \item En \textbf{Body $>$ raw $>$ JSON} pegue:
\end{enumerate}

\begin{lstlisting}
{
  "periodo_inicio": "2026-12-31",
  "periodo_fin": "2026-01-01"
}
\end{lstlisting}

\esperado{Status \textbf{400}. La fecha final no puede ser anterior a la inicial.}

\subsection{META-01 --- Crear una meta vinculada al objetivo}

\begin{enumerate}
  \item Cree una solicitud \textbf{POST} a \url{http://localhost:8000/api/metas/}.
  \item En \textbf{Headers} agregue \path{Content-Type = application/json} y \path{X-CSRFToken = token actual}.
  \item En \textbf{Body $>$ raw $>$ JSON} pegue:
\end{enumerate}

\begin{lstlisting}
{
  "plan": <ID_PLAN>,
  "objetivo_estrategico": <ID_OBJETIVO_CREADO>,
  "nombre": "Mejorar procesos institucionales GUIA-01",
  "descripcion": "Meta de la guia de pruebas.",
  "resultado_esperado": "Obtener una mejora medible durante 2026.",
  "fecha_inicio": "2026-01-01",
  "fecha_fin": "2026-12-31"
}
\end{lstlisting}

\esperado{Status \textbf{201}, estado BORRADOR y detalle del objetivo estratégico relacionado. Copie el ID de la meta.}

\subsection{META-02 --- Activar la meta}

\begin{enumerate}
  \item Cree una solicitud \textbf{POST} a \path{http://localhost:8000/api/metas/<ID_META>/activar/}.
  \item En \textbf{Headers} agregue \path{Content-Type = application/json} y \path{X-CSRFToken = token actual}.
  \item En \textbf{Body $>$ raw $>$ JSON} escriba \texttt{\{\}}.
  \item Pulse \textbf{Send}.
\end{enumerate}

\esperado{Status \textbf{200}, estado ACTIVA y \path{activa: true}.}

\subsection{IND-01 --- Crear un indicador}

\begin{enumerate}
  \item Cree una solicitud \textbf{POST} a \url{http://localhost:8000/api/indicadores/}.
  \item En \textbf{Headers} agregue \path{Content-Type = application/json} y \path{X-CSRFToken = token actual}.
  \item En \textbf{Body $>$ raw $>$ JSON} pegue:
\end{enumerate}

\begin{lstlisting}
{
  "meta": <ID_META>,
  "nombre": "Porcentaje de procesos mejorados GUIA-01",
  "descripcion": "Indicador de la meta creada en Postman.",
  "unidad_medida": "Porcentaje",
  "valor_base": "0.00",
  "valor_meta": "100.00",
  "frecuencia": "MENSUAL"
}
\end{lstlisting}

\esperado{Status \textbf{201}. Copie el ID del indicador. Su \path{valor_actual} debe ser 0.00.}

\subsection{AVA-01 --- Rechazar un avance negativo}

\begin{enumerate}
  \item Cree una solicitud \textbf{POST} a \path{http://localhost:8000/api/indicadores/<ID_INDICADOR>/registrar-avance/}.
  \item En \textbf{Headers} agregue \path{Content-Type = application/json} y \path{X-CSRFToken = token actual}.
  \item En \textbf{Body $>$ raw $>$ JSON} pegue:
\end{enumerate}

\begin{lstlisting}
{
  "fecha_registro": "2026-07-15",
  "valor": "-1.00",
  "observacion": "Valor negativo enviado intencionalmente."
}
\end{lstlisting}

\esperado{Status \textbf{400}. No debe crearse el avance.}

\subsection{AVA-02 --- Registrar un avance correcto}

Mantenga el método, URL y headers anteriores. Cambie el Body por:

\begin{lstlisting}
{
  "fecha_registro": "2026-07-15",
  "valor": "25.00",
  "observacion": "Primer avance registrado desde Postman."
}
\end{lstlisting}

\esperado{Status \textbf{201}. Compruebe \path{registrado_por_detalle.username: "planificador"}; el registrador no se elige desde el Body.}

Ahora cree una solicitud \textbf{GET}. No seleccione Body ni agregue filas manuales en \textbf{Headers}. Use esta URL:

\begin{lstlisting}
http://localhost:8000/api/indicadores/<ID_INDICADOR>/
\end{lstlisting}

\esperado{Status \textbf{200} y \path{valor_actual: "25.00"}. Esto comprueba que el avance actualizó el indicador.}

\subsection{META-03 --- Intentar eliminar una meta que ya tiene indicador}

El perfil Planificador no tiene permiso de eliminación. Para probar específicamente la protección de relaciones y obtener el conflicto 409, cambie temporalmente a \path{superadministrador}, haga el login y copie su nuevo token CSRF.

\begin{enumerate}
  \item Cree una solicitud \textbf{DELETE} a \path{http://localhost:8000/api/metas/<ID_META>/}.
  \item En \textbf{Headers} agregue \path{Content-Type = application/json} y \path{X-CSRFToken = token actual}.
  \item No agregue Body.
  \item Pulse \textbf{Send}.
\end{enumerate}

\esperado{Status \textbf{409}. La meta y el indicador se conservan porque ya existe información relacionada.}

\subsection{PLAN-03 --- Enviar el plan a revisión}

Vuelva a cambiar al usuario \path{planificador} y copie el token CSRF obtenido después de ese login.

\begin{enumerate}
  \item Cree una solicitud \textbf{POST} a \path{http://localhost:8000/api/planes/<ID_PLAN>/enviar-a-revision/}.
  \item En \textbf{Headers} agregue \path{Content-Type = application/json} y \path{X-CSRFToken = token actual}.
  \item En \textbf{Body $>$ raw $>$ JSON} pegue:
\end{enumerate}

\begin{lstlisting}
{
  "observacion": "Plan enviado a revision desde Postman."
}
\end{lstlisting}

\esperado{Status \textbf{200} y estado EN\_REVISION.}

Pulse \textbf{Send} una segunda vez sin cambiar nada.

\esperado{Status \textbf{409}, porque un plan que ya está en revisión no puede enviarse nuevamente.}

Finalmente cambie la URL a \path{http://localhost:8000/api/planes/<ID_PLAN>/aprobar/}, mantenga POST, los dos headers y use Body \texttt{\{\}}.

\esperado{Status \textbf{403}. El Planificador puede enviar, pero no aprobar.}

\section{Prueba 8: revisión y aprobación con Supervisor}

Cambie al usuario \path{supervisor}. Esta prueba usa el plan de SNP que inicialmente se encuentra en EN\_REVISION.

\begin{enumerate}
  \item Cree una solicitud \textbf{GET} a:
\end{enumerate}

\begin{lstlisting}
http://localhost:8000/api/planes/?search=Plan%20Institucional%202026-2029
\end{lstlisting}

En la respuesta elija el plan cuya entidad sea SNP y cuyo estado sea EN\_REVISION. Copie su ID.

Para iniciar la revisión:

\begin{enumerate}
  \item Cree una solicitud \textbf{POST} a \path{http://localhost:8000/api/planes/<ID_PLAN_SNP>/revisar/}.
  \item En \textbf{Headers} agregue \path{Content-Type = application/json} y \path{X-CSRFToken = token actual}.
  \item En \textbf{Body $>$ raw $>$ JSON} escriba \texttt{\{\}}.
  \item Pulse \textbf{Send}.
\end{enumerate}

\esperado{Status \textbf{200} y estado EN\_REVISION\_INICIADA.}

Para aprobarlo, cambie únicamente la URL a:

\begin{lstlisting}
http://localhost:8000/api/planes/<ID_PLAN_SNP>/aprobar/
\end{lstlisting}

Mantenga POST, los dos headers y el Body \texttt{\{\}}.

\esperado{Status \textbf{200} y estado APROBADO.}

Para comprobar el historial cree una solicitud \textbf{GET}. No seleccione Body ni agregue filas manuales en \textbf{Headers}. Use esta URL:

\begin{lstlisting}
http://localhost:8000/api/planes/<ID_PLAN_SNP>/historial/
\end{lstlisting}

\esperado{Aparecen las acciones de iniciar revisión y aprobar, realizadas por \path{supervisor}.}

Si el plan ya estaba aprobado antes de empezar, no fuerce su estado con PATCH. Use otro plan de SNP que haya llegado correctamente a EN\_REVISION o deje esta prueba anotada como ya ejecutada.

\section{Prueba 9: seguimiento de un proyecto}

Cambie al usuario \path{planificador}. Use los IDs de PIR-001 y su hito que obtuvo en la Prueba 2.

Primero consulte el proyecto mediante una solicitud \textbf{GET}. No seleccione Body ni agregue filas manuales en \textbf{Headers}:

\begin{lstlisting}
http://localhost:8000/api/proyectos/?search=PIR-001
\end{lstlisting}

\esperado{Status \textbf{200}. El proyecto aparece en estado EN\_EJECUCION.}

Si PIR-001 ya no está EN\_EJECUCION porque otra prueba lo suspendió, finalizó o archivó, no continúe con ese registro: utilice otro proyecto en ejecución o vuelva a preparar los datos demostrativos locales.

Para registrar un seguimiento:

\begin{enumerate}
  \item Cree una solicitud \textbf{POST} a \path{http://localhost:8000/api/proyectos/<ID_PIR_001>/registrar-seguimiento/}.
  \item En \textbf{Headers} agregue \path{Content-Type = application/json} y \path{X-CSRFToken = token actual}.
  \item En \textbf{Body $>$ raw $>$ JSON} pegue:
\end{enumerate}

\begin{lstlisting}
{
  "hito": <ID_HITO_PIR_001>,
  "fecha_registro": "2026-07-20",
  "avance_fisico": "40.00",
  "avance_financiero": "37.00",
  "observacion": "Corte registrado desde Postman."
}
\end{lstlisting}

\esperado{Status \textbf{201}. El usuario registrador debe proceder de la sesión. Al consultar nuevamente el proyecto, sus avances consolidados deben ser 40.00 y 37.00.}

Pulse \textbf{Send} nuevamente con la misma fecha.

\esperado{Status \textbf{400}, porque solo puede existir un seguimiento por proyecto y fecha.}

Si el 20 de julio ya está registrado, use otra fecha de 2026 que esté dentro del periodo del proyecto, no sea futura y no se haya utilizado antes.

\section{Prueba 10: reportes y protección entre instituciones}

Cambie al usuario \path{externo}.

\subsection{REP-01 --- Consultar los reportes disponibles}

Cree una solicitud \textbf{GET}. No seleccione Body ni agregue filas manuales en \textbf{Headers}. Use esta URL:

\begin{lstlisting}
http://localhost:8000/api/reportes/catalogo/
\end{lstlisting}

\esperado{Status \textbf{200}. Solo aparecen los reportes permitidos para el Usuario Externo.}

\subsection{REP-02 --- Ver planes aprobados de su institución}

Cree una solicitud \textbf{GET}. No seleccione Body ni agregue filas manuales en \textbf{Headers}. Use esta URL:

\begin{lstlisting}
http://localhost:8000/api/reportes/generar/planes/?estado=APROBADO&limite=50
\end{lstlisting}

\esperado{Status \textbf{200}. Los resultados pertenecen a MDI y no incluyen planes de SNP.}

\subsection{REP-03 --- Intentar ampliar el alcance mediante un parámetro}

Use una solicitud \textbf{GET} a:

\begin{lstlisting}
http://localhost:8000/api/reportes/generar/planes/?entidad=<ID_ENTIDAD_SNP>&limite=50
\end{lstlisting}

\esperado{Status \textbf{200}, \path{total: 0} y resultados vacíos. Escribir el ID de SNP en la URL no concede acceso a SNP.}

\subsection{REP-04 --- Descargar un reporte}

Use una solicitud \textbf{GET} a:

\begin{lstlisting}
http://localhost:8000/api/reportes/exportar/planes/csv/?entidad=<ID_ENTIDAD_MDI>&limite=50
\end{lstlisting}

Para guardar el archivo, abra la flecha situada junto al botón \textbf{Send}, seleccione \textbf{Send and Download} y elija la ubicación. También puede usar \textbf{Send} y después \textbf{Save Response}.

\esperado{Status \textbf{200} y descarga de un archivo CSV con datos autorizados de MDI.}

Para probar los otros formatos, cambie únicamente \path{csv} por \path{json}, \path{xlsx} o \path{pdf}.

\section{Prueba 11: auditoría y perfil de solo lectura}

Cambie al usuario \path{auditor}.

\subsection{AUD-01 --- Consultar eventos de autenticación}

Cree una solicitud \textbf{GET}. No seleccione Body ni agregue filas manuales en \textbf{Headers}. Use esta URL:

\begin{lstlisting}
http://localhost:8000/api/auditoria/eventos/?modulo=autenticacion&ordering=-fecha_hora
\end{lstlisting}

\esperado{Status \textbf{200}. Aparecen eventos de autenticación que están dentro del alcance SNP. Ningún evento debe contener contraseñas, cookies o tokens.}

\subsection{AUD-02 --- Comprobar que el Auditor no puede crear planes}

\begin{enumerate}
  \item Cree una solicitud \textbf{POST} a \url{http://localhost:8000/api/planes/}.
  \item En \textbf{Headers} agregue \path{Content-Type = application/json} y \path{X-CSRFToken = token actual}.
  \item En \textbf{Body $>$ raw $>$ JSON} pegue:
\end{enumerate}

\begin{lstlisting}
{
  "nombre": "Plan que no debe crear el auditor",
  "descripcion": "Intento controlado.",
  "entidad": <ID_ENTIDAD_SNP>,
  "periodo_inicio": "2026-01-01",
  "periodo_fin": "2026-12-31",
  "responsable": <ID_USUARIO_SUPERVISOR>
}
\end{lstlisting}

\esperado{Status \textbf{403}. El Auditor puede leer información autorizada, pero no crear ni modificar datos.}

\section{Qué significa cada status que aparecerá}

\begin{center}
\begin{tabularx}{\textwidth}{p{1.7cm} Y}
\toprule
\textbf{Status} & \textbf{Interpretación sencilla} \\
\midrule
200 & La consulta o acción terminó correctamente. \\
201 & Se creó un registro correctamente. \\
204 & Se eliminó el registro; es normal que la respuesta no tenga Body. \\
400 & El JSON contiene un dato inválido o no cumple una regla del negocio. Lea el mensaje del campo. \\
401 & No existe una sesión válida. Debe iniciar sesión. \\
403 & El usuario no tiene permiso o el token CSRF es incorrecto. Revise el campo \path{code}. \\
404 & El ID no existe o el registro está fuera del alcance institucional del usuario. \\
409 & La acción entra en conflicto con el estado actual o con relaciones que deben conservarse. \\
429 & Se hicieron demasiados intentos de login; debe esperar antes de repetir. \\
\bottomrule
\end{tabularx}
\end{center}

\section{Si una solicitud no responde como esperaba}

\begin{itemize}
  \item Si recibe 403 con \path{csrf_failed}, ejecute otra vez \path{GET /api/auth/csrf/}, copie el token nuevo y reemplace \path{X-CSRFToken}.
  \item Si recibe 401 después de un login correcto, revise que todas las URLs utilicen \path{localhost} y que Postman conserve la cookie \path{sessionid}.
  \item Si recibe 404 con un ID que sabe que existe, compruebe si el usuario pertenece a otra institución.
  \item Si recibe 400, confirme que reemplazó todos los marcadores \path{<ID_...>} por números reales y que seleccionó \textbf{Body $>$ raw $>$ JSON}.
  \item Si recibe 409 al repetir una prueba, consulte primero el estado actual. La primera solicitud probablemente ya cambió el registro.
  \item Si un código o nombre \path{GUIA-01} ya existe, cambie el sufijo a \path{GUIA-02} en todos los registros relacionados de esa nueva ejecución.
\end{itemize}

\vfill
\begin{center}
\textcolor{AzulSipeip}{\rule{0.65\textwidth}{0.5pt}}\\
\small Fin de la guía práctica
\end{center}

\end{document}
